\documentclass[ngerman,xcolor={table,dvipsnames},smaller,compress,hyperref={bookmarks,colorlinks}]{beamer}%,handout

\usepackage{listings}
\usepackage[utf8]{inputenc}
\usepackage[english]{babel}
\usepackage{mathpazo}
\usepackage{textcomp}
\usepackage{enumerate}
\usepackage{textgreek}
\usepackage{xfrac}
%\usepackage[hyphens]{url}

%\usepackage{auto-pst-pdf}
\usepackage{graphicx}
\usepackage{pstricks}
\usepackage{pst-node}
\usepackage{pst-uml}
\usepackage{pst-tree}
\usepackage{pst-blur}

\usepackage{amsfonts}
\usepackage{amssymb}
\usepackage{amsmath}
\usepackage{pifont}

% for multipage tables (xtab or longtable
\usepackage{longtable}
\usepackage{lscape}
\usepackage{booktabs}
\usepackage[tableposition=top]{caption}

\usepackage{hyperref}
\usepackage[hyphenbreaks]{breakurl}

\lstdefinelanguage{Perl6}%
  {morekeywords={eager,hyper,race,my,say,dd},%
   sensitive,%
   morecomment=[l]\#,%
   morestring=[b]",%
   morestring=[b]',%
   MoreSelectCharTable=%
      \lst@ReplaceInput{\$\#}{\lst@ProcessOther\$\lst@ProcessOther\#}%
  }[keywords,comments,strings]%

\title{Perl6 for Beginners}
\author{Jens Rehsack}
\institute[Niederrhein.PM]{Niederrhein Perl Mongers}
\date{European Perl Conference 2019}

\usetheme[secheader]{Boadilla}
\setbeamertemplate{navigation symbols}{}

\newcommand{\perlfilename}[1]{{\color{cyan}{\textit{\begingroup \urlstyle{sf}\Url{#1}}}}}
\newcommand{\xsfilename}[1]{{\color{olive}{\textit{\begingroup \urlstyle{sf}\Url{#1}}}}}
\newcommand{\hfilename}[1]{{\color{olive}{\textit{\begingroup \urlstyle{sf}\Url{#1}}}}}
\newcommand{\cfilename}[1]{{\color{magenta}{\textit{\begingroup \urlstyle{sf}\Url{#1}}}}}
\newcommand{\variable}[1]{{\color{violet}{\textsf{#1}}}}
\newcommand{\variablew}[1]{{{\textsf{#1}}}}
\newcommand{\command}[1]{'\texttt{#1}'}

\def\signed #1{{\leavevmode\unskip\nobreak\hfil\penalty50\hskip2em
   \hbox{}\nobreak\hfil(#1)%
   \parfillskip=0pt \finalhyphendemerits=0 \endgraf}}
 
\newsavebox\mybox
\newenvironment{aquote}[1]
   {\savebox\mybox{#1}\begin{quotation}}
   {\signed{\usebox\mybox}\end{quotation}}

\makeatletter
\makeatother

\begin{document}

\psset{angleA=90,angleB=-90}
\lstset{language=Perl,
        basicstyle=\ttfamily,
        keywordstyle=\color{Maroon},
        commentstyle=\color{Blue}, 
        stringstyle=\color{Green},
        showstringspaces=false}

\lstset{
    literate=%
        {Ö}{{\"O}}1
        {Ä}{{\"A}}1
        {Ü}{{\"U}}1
        {ß}{{\ss}}2
        {ü}{{\"u}}1
        {ä}{{\"a}}1
        {ö}{{\"o}}1
        {°}{{\textdegree}}1
        {α}{{\textalpha}}1
        {β}{{\textbeta}}1
        {γ}{{\textgamma}}1
        {δ}{{\textdelta}}1
        {ε}{{\textepsilon}}1
        {ζ}{{\textzeta}}1
        {η}{{\texteta}}1
        {θ}{{\texttheta}}1
        {ι}{{\textiota}}1
        {κ}{{\textkappa}}1
        {λ}{{\textlambda}}1
        {μ}{{\textmugreek}}1
        {ν}{{\textnu}}1
        {ξ}{{\textxi}}1
        {ο}{{\textomikron}}1
        {π}{{\textpi}}1
        {ρ}{{\textrho}}1
        {σ}{{\textsigma }}1
        {τ}{{\texttau}}1
        {υ}{{\textupsilon}}1
        {φ}{{\textphi}}1
        {χ}{{\textchi}}1
        {ψ}{{\textpsi}}1
        {ω}{{\textomega}}1
        {Α}{{\textAlpha}}1
        {Β}{{\textBeta}}1
        {Γ}{{\textGamma}}1
        {Δ}{{\textDelta}}1
        {Ε}{{\textEpsilon}}1
        {Ζ}{{\textZeta}}1
        {Η}{{\textEta}}1
        {Θ}{{\textTheta}}1
        {Ι}{{\textIota}}1
        {Κ}{{\textKappa}}1
        {Λ}{{\textLambda}}1
        {Μ}{{\textMu}}1
        {Ν}{{\textNu}}1
        {Ξ}{{\textXi}}1
        {Ο}{{\textOmikron}}1
        {Π}{{\textPi}}1
        {Ρ}{{\textRho}}1
        {Σ}{{\textSigma}}1
        {Τ}{{\textTau}}1
        {Υ}{{\textUpsilon}}1
        {Φ}{{\textPhi}}1
        {Χ}{{\textChi}}1
        {Ψ}{{\textPsi}}1
        {Ω}{{\textOmega}}1
        {ς}{{\textvarsigma}}1
        {φ}{{\straightphi}}1
        {θ}{{\scripttheta}}1
        {θ}{{\straighttheta}}1
        {ε}{{\straightepsilon}}1
        {¹}{{\textsuperscript{1}}}1
        {²}{{\textsuperscript{2}}}1
        {³}{{\textsuperscript{3}}}1
        {⁴}{{\textsuperscript{4}}}1
        {⁵}{{\textsuperscript{5}}}1
        {⁶}{{\textsuperscript{6}}}1
        {⁷}{{\textsuperscript{7}}}1
        {⁸}{{\textsuperscript{8}}}1
        {⁹}{{\textsuperscript{9}}}1
        {⁰}{{\textsuperscript{0}}}1
        {½}{{\sfrac{1}{2}}}1
}


\AtBeginPart{\begin{frame}<beamer> \frametitle{Overview} \partpage \tableofcontents[current] \end{frame}}

\frame{\maketitle}

\part{Introduction}

\section{Introduction}

\subsection{Motivation}

\begin{frame}[fragile]{Motivation}
\begin{block}<1->{Collecting (Lost) People}
\begin{itemize}
\item<1-> Perl 6 is in development for more than 10 years and it has a lot of
          impressive new features.
\item<2-> The language syntax evolved with the language capabilities and the
          evolved mindset.
\item<3-> Getting snippets of code from experienced Perl 6 developers might create
          more (new) questions than answers (if any).
\item<4-> This talk attempts to start at the beginning and collects everyone
          on the way.
\end{itemize}
\end{block}

\begin{block}<5->{Goals}
At the end of the talk, attendees should be able to write small Perl 6 code
and understand given examples.
\end{block}
\end{frame}

\subsection{Is \& Isn't}

\begin{frame}[fragile]{What Perl 6 Is}
\begin{block}<1->{Perl 6 Is \ldots}
\begin{itemize}
\item<1-> a modern language of the Perl family
\item<2-> a language with a fast growing module ecosystem
\item<3-> a development environment with built-in support for asynchronouos I/O,
          auto parallelization, multiple syntaxes in one program, \ldots
\item<4-> a name of a open minded, strong collaborating community
\end{itemize}
\end{block}
\end{frame}

\begin{frame}[fragile]{What Perl (6) Is not}
\begin{block}<1->{Perl 6 Isn't \ldots}
\begin{itemize}
\item<1-> Perl 4
\item<2-> Perl 5
\item<3-> intended to be a drop-in replacement of any of them
\item<4-> continuity of ''Status Quo''
\end{itemize}
\end{block}
\end{frame}

\begin{frame}[fragile]{What Rakudo Is}
\begin{block}<1->{Rakudo Is \ldots}
\begin{itemize}
\item<1-> a Perl 6 compiler (targeting Moar VM and JVM)
\item<2-> Perl 6 Specification conform
\item<3-> one of many Perl 6 implementations
\end{itemize}
\end{block}

\begin{block}<4->{Rakudo* Is \ldots}
\begin{itemize}
\item<4-> a bundle of the Perl 6 compiler and curated Perl 6 ''core'' modules
\item<5-> what you might call ''perl''
\end{itemize}
\end{block}
\end{frame}

\section{Starting}

\begin{frame}[fragile]{Getting \& Installation (Binaries)}
\begin{block}<1->{macOS \& Windows}
\begin{itemize}
\item<1-> Go to \url{https://rakudo.org/files}
\item<2-> Choose your Operating System or Windows
\item<3-> Follow the installation instructions for your Operating System or Windows
\end{itemize}
\end{block}
\end{frame}

%\begin{frame}[fragile]{Packages Repository I/II}
%\begin{block}<1->{*BSD, Linux, \ldots - use package repository}
%\begin{lstlisting}[language=sh,inputencoding=latin9]
%nbsd   $ export PKG_PATH="http://cdn.NetBSD.org/pub/\
%         pkgsrc/packages/OPSYS/ARCH/VERSIONS/All/"
%fbsd   $ # pkg knows it's repository
%obsd   $ export PKG_PATH=scp://user@company-build-\
%         server/usr/ports/packages/\%a/all:\
%	 https://trusted-public-server/\%m:\
%	 installpath
%debian $ sudo apt-get -y update
%centos $ # yum updates repositories automatically
%suse   $ sudo zypper refresh
%\end{lstlisting}
%\end{block}
%\end{frame}
%
%\begin{frame}[fragile]{Packages Repository II/II}
%\begin{block}<1->{*BSD, Linux, \ldots - install from package repository}
%\begin{lstlisting}[language=sh,inputencoding=latin9]
%nbsd   $ sudo pkg_add rakudo-star
%fbsd   $ sudo pkg install rakudo-star
%obsd   $ sudo pkg_add rakudo-star
%debian $ sudo apt-get install rakudo-star
%centos $ sudo yum install rakudo star
%suse   $ sudo zypper install rakudo-star
%\end{lstlisting}
%\end{block}
%\end{frame}
%
%\begin{frame}[fragile]{Install from Source}
%\begin{block}<1->{Install Rakudo Star from Source}
%There is an excellent guide at \url{https://rakudo.org/files/star/source}
%for most widely used operating system distributions. \uncover<2->{Since
%there is the source download included, an offline copy doesn't make much
%sense.}
%
%\uncover<3->{\textrightarrow DVD distributors might ask for an offline copy
%for the targeted audience.}
%\end{block}
%\end{frame}
%
%\begin{frame}[fragile]{Manage multiple Rakudo Installations}
%\begin{block}<1->{Brewing continues \ldots}
%\begin{itemize}
%\item<1-> Visit \url{https://github.com/tadzik/rakudobrew}
%\item<2-> Install \texttt{git} unless already available
%\item<3-> run \lstinline[language=sh,inputencoding=latin9]!git clone https://github.com/tadzik/rakudobrew ~/.rakudobrew!
%(Windows users might prefer \lstinline[language=sh,inputencoding=latin9]!%USERPROFILE%\rakudobrew!
%instead of \lstinline[language=sh,inputencoding=latin9]!~/.rakudobrew! in this \uncover<4->{and next command})
%\item<4-> run ~/.rakudobrew/bin/init! as described in \texttt{README.md}, e.g.
%\lstinline[language=sh,inputencoding=latin9]!echo 'eval "$(~/.rakudobrew/bin/rakudobrew init Bash)"' \!
%\lstinline[language=sh,inputencoding=latin9]!  >> ~/.profile!
%\end{itemize}
%\end{block}
%\end{frame}



\section{Example Code}

\subsection{Running Perl 6}

\begin{frame}[fragile]{Starting Perl 6 Program}
\begin{block}<1->{How to run Perl 6 Programs}
\begin{itemize}
\item<1-> Run explicitely with compiler:
\begin{lstlisting}[language=sh,inputencoding=latin9]
sno@ernie$ perl6 prog.p6
\end{lstlisting}
\item<2-> Run self-contained
\begin{lstlisting}[language=sh,inputencoding=latin9]
sno@ernie$ head -n1 ./prog.p6
#!/opt/pkg/bin/perl6
sno@ernie$ chmod +x ./prog.p6
sno@ernie$ ./prog.p6
\end{lstlisting}
\end{itemize}
\end{block}
\end{frame}

\begin{frame}[fragile]{Starting Program With Parameters}
\begin{block}<1->{How to run Perl 6 Programs With Parameters}
\begin{itemize}
\item<1-> Run explicitely with compiler:
\begin{lstlisting}[language=sh,inputencoding=latin9]
sno@ernie$ perl6 prog.p6 --arg1 --arg2=val pos_arg
\end{lstlisting}
\item<2-> Run self-contained
\begin{lstlisting}[language=sh,inputencoding=latin9]
sno@ernie$ head -n1 ./prog.p6
#!/opt/pkg/bin/perl6
sno@ernie$ chmod +x ./prog.p6
sno@ernie$ ./prog.p6 --arg1 --arg2=val pos_arg
\end{lstlisting}
\end{itemize}
\end{block}
\end{frame}

\begin{frame}[fragile]{Starting Perl6 with Parameters}
\begin{block}<1->{Starting Perl6 with Parameters}
\begin{itemize}
\item<1-> Perl 6 specifies in \url{https://github.com/perl6/specs/blob/master/S19-commandline.pod}
which parameters are available and what they do.
\item<2-> A more detailed \burlalt{https://github.com/perl6/specs/blob/master/S19-commandline.pod#Option-Reference}{Option Reference} is available later on that page - so scroll down and read until the end!
\item<3-> However, Rakudo implements some options (\texttt{-M}) and refuses some options (\texttt{-F}) - so always refer to \lstinline[language=sh,inputencoding=latin9]!perl6 --help!
\item<4-> Change perspective: Some things are better solved differently than adding options \lstinline[language=sh,inputencoding=latin9]!perl6-debug ./prog.p6! replaces \lstinline[language=sh,inputencoding=latin9]!perl -d ./prog.pl!
\end{itemize}
\end{block}
\end{frame}


\subsection{One Liners}

\begin{frame}[fragile]{Killer App: One Liners}
\begin{block}<1->{Doing One Liners with Perl 6}
Perl has always been famous for it's one-liners. Let's continue \ldots
\begin{itemize}
\item<1-> Strip trailing whitespaces \lstinline[language=sh,inputencoding=latin9]!perl6 -pe 's/\s*$//'!
Looks familiar, doesn't it?
\item<2-> Most use is very likely the line ending tidying - done easily in Perl 6 by \lstinline[language=sh,inputencoding=latin9]!perl6 -pe ''!
Yes - it's completely built-in. An empty by-line job does \lstinline[language=sh,inputencoding=latin9]!perl5  -p -i -e 's/\012?\015/\n/g'!
\item<3-> Another famous job is filtering, like \lstinline[language=sh,inputencoding=latin9]!perl6 -ne '.say if .is-prime'!
\item<4-> Adding line numbers to some inputfile needed? How about
\begin{lstlisting}[language=sh,inputencoding=latin9]
perl6 -ne 'for lines.kv -> $no, $l {\
say sprintf("%04d: %s", $no, $l) }'
\end{lstlisting}
\item<5-> Maybe a bit more flexibility in width of line-numbers?
\begin{lstlisting}[language=sh,inputencoding=latin9]
perl6 -e 'my @l = lines; \
my $fmt = "\%0" ~ @l.elems.chars ~ "d: %s"; \
for ^@l.elems {say sprintf($fmt, $_+1, @l[$_]) }'
\end{lstlisting}
\end{itemize}
\end{block}
\end{frame}



\subsection{Strings}

\begin{frame}[fragile]{Defining Strings}
\begin{block}<1->{Strings -- Synopsis}
\small
\begin{lstlisting}[language=Perl6,inputencoding=utf8,escapeinside={(*@}{@*)}]
dd (*@\pnode(0,0){A1}{}@*)'I can see $(*@\pnode(0,0){A2}{}@*)$ in your eyes';
dd (*@\pnode(0,0){B}{}@*)"I am N° $*(*@\pnode(0,0){C}{}@*)PID";
my $s = "Haa"; $s(*@\pnode(0,0){D}{}@*)++; say "Interpolated <$(*@\pnode(0,0){E}{}@*)s>!";

dd q(*@\pnode(0,0){F}{}@*):to/TERMINATOR(*@\pnode(0,0){G}{}@*)/
Lorem ipsum dolor sit amet, consectetur adipisici elit,
sed eiusmod tempor incidunt ut labore et dolore magna
aliqua.
TERMINATOR
\end{lstlisting}
\end{block}

\begin{itemize}
\item<2-> This \rnode{a1}{encompasses} a literal string -- the string will be processed \rnode{a2}{exactly} as written.
\item<3-> That string \rnode{b}{encompassing} causes interpolation. \uncover<4->{Mind the \rnode{c}{\texttt{\$*PID}} instead of the \texttt{\$\$} from Perl 5 or Bourne Shell constructs.}
\item<5->  \rnode{d}{Increments} \texttt{Haa} \textrightarrow{}  \texttt{Hab} and prints \textit{Interpolated  \rnode{e}{\textless{}Hab\textgreater{}}!}
\item<6-> \rnode{g}{Heredocs} are included into \texttt{\rnode{f}{q}} and \texttt{qq}, respectively
\end{itemize}
\uncover<2>{\nccurve[linecolor=teal]{->}{a1}{A1}\nccurve[linecolor=teal]{->}{a2}{A2}}
\uncover<3-4>{\nccurve[linecolor=teal]{->}{b}{B}}
\uncover<4>{\nccurve[linecolor=teal]{->}{c}{C}}
\uncover<5>{\nccurve[linecolor=teal]{->}{d}{D}\nccurve[linecolor=teal]{->}{e}{E}}
\uncover<6>{\nccurve[linecolor=teal]{->}{f}{F}\nccurve[linecolor=teal]{->}{g}{G}}
\end{frame}

\begin{frame}[fragile]{C special characters \& substr}
\begin{block}<1->{Strings -- Synopsis}
\small
\begin{lstlisting}[language=Perl6,inputencoding=utf8,escapeinside={(*@}{@*)}]
my $s = "(*@\pnode(0,0){A1}{}@*)\a, \b, \e, \f, \n, \r, \t, \\, \x1b, \cJ(*@\pnode(0,0){A2}{}@*)";

my $bell = substr($s, (*@\pnode(0,0){B}{}@*)0, 1);     dd $bell;
my $nl   = substr($s, (*@\pnode(0,0){C}{}@*)*-1, 1);   dd $nl;
my $tab  = substr($s, (*@\pnode(0,0){D}{}@*)18..18);   dd $tab;
my $lf   = substr($s, 15, (*@\pnode(0,0){E}{}@*)*-12); dd $lf;
\end{lstlisting}
\end{block}

\begin{itemize}
\item<2-> \ldots well known \rnode{a}{list of special characters} for terminal controlling.
\item<3-> Getting the \rnode{b}{first} character (bell)
\item<4-> Getting the \rnode{c}{last} character (newline)\uncover<5->(how many newlines are in that string?)
\item<6-> Fetching a \rnode{d}{range} of 1 character containing a tabulator special char from the middle
\item<7-> Beware \rnode{e}{negative} addressing of \texttt{to} --- getting line feed character
\end{itemize}
\uncover<2>{\nccurve[linecolor=teal]{->}{a}{A1}\nccurve[linecolor=teal]{->}{a}{A2}}
\uncover<3>{\nccurve[linecolor=teal]{->}{b}{B}}
\uncover<4-5>{\nccurve[linecolor=teal]{->}{c}{C}}
\uncover<6>{\nccurve[linecolor=teal]{->}{d}{D}}
\uncover<7>{\nccurve[linecolor=teal]{->}{e}{E}}
\end{frame}

\begin{frame}[fragile]{''Attributes'' of Strings}
\begin{block}<1->{Strings -- Synopsis}
\small
\begin{lstlisting}[language=Perl6,inputencoding=utf8,escapeinside={(*@}{@*)}]
my $s = "\c[WAVING WHITE FLAG, VARIATION SELECTOR-16,
            ZERO WIDTH JOINER, RAINBOW]";
say $s;
say "chars: ", (*@\pnode(0,0){A}{}@*)$s.chars;
say "codes: ", (*@\pnode(0,0){B}{}@*)$s.codes;
say "encode.bytes: ", (*@\pnode(0,0){C}{}@*)$s.encode.bytes;
say "encode('utf-16').bytes: ", (*@\pnode(0,0){D}{}@*)$s.encode('utf-16').bytes;
\end{lstlisting}
\end{block}

\ldots on a typical unicode character \uncover<2->{(\includegraphics[height=\baselineskip]{rainbow-flag-apple.eps})}.
\begin{itemize}
\item<3-> It is \texttt{\rnode{a}{1}} character
\item<4-> containing \texttt{\rnode{b}{4}} codes
\item<5-> using \texttt{\rnode{c}{14}} bytes in default (UTF-8) encoding
\item<6-> or \texttt{\rnode{d}{12}} bytes in UTF-16 encoding.
\end{itemize}
\uncover<3>{\nccurve[linecolor=teal]{->}{a}{A}}
\uncover<4>{\nccurve[linecolor=teal]{->}{b}{B}}
\uncover<5>{\nccurve[linecolor=teal]{->}{c}{C}}
\uncover<6>{\nccurve[linecolor=teal]{->}{d}{D}}
\end{frame}

\begin{frame}[fragile]{String quoting}
\begin{block}<1->{Strings -- Synopsis}
\small
\begin{lstlisting}[language=Perl6,inputencoding=utf8,escapeinside={(*@}{@*)}]
my $l = (*@\pnode(0,0){A}{}@*)Q (SELECT COUNT(*@\pnode(0,0){B}{}@*)(*) FROM users);
my @f = (*@\pnode(0,0){C}{}@*)Q:(*@\pnode(0,0){D}{}@*)x:w(*@\pnode(0,0){E}{}@*)w/l(*@\pnode(0,0){F}{}@*)s/;
my $e = (*@\pnode(0,0){G}{}@*)q<swept volume (*@\pnode(0,0){H}{}@*)\>(*@\pnode(0,0){I}{}@*) turbocharger>
my $i = (*@\pnode(0,0){J}{}@*)qq:!(*@\pnode(0,0){k}{}@*)h[%(*@\pnode(0,0){l}{}@*)0{@f.elems.chars}d: %s]
\end{lstlisting}
\end{block}

\begin{itemize}
\item<2-> \rnode{a}{Literal string} -- with \rnode{b}{nested delimiters}.
\item<3-> \rnode{c}{Literal} \rnode{f}{string} \textrightarrow{} \rnode{d}{executed} \textrightarrow{} \rnode{e}{split on words (with quote protection)}
\item<4-> \rnode{g}{Almost literal} string, but allows \rnode{h}{escaping} (e.g. to ''fix'' \rnode{i}{delimiters used} in string)
\item<5-> Fully \rnode{j}{interpolated} string -- \rnode{k}{except hashes} for this \rnode{l}{special} case.
\end{itemize}
\uncover<2>{\nccurve[linecolor=teal]{->}{a}{A}\nccurve[linecolor=teal]{->}{b}{B}}
\uncover<3>{\nccurve[linecolor=teal]{->}{c}{C}\nccurve[linecolor=violet]{->}{d}{D}\nccurve[linecolor=violet]{->}{e}{E}\nccurve[linecolor=teal]{->}{f}{F}}
\uncover<4>{\nccurve[linecolor=teal]{->}{g}{G}\nccurve[linecolor=violet]{->}{h}{H}\nccurve[linecolor=violet]{->}{i}{I}}
\uncover<5>{\nccurve[linecolor=teal]{->}{j}{J}\nccurve[linecolor=violet]{->}{k}{K}\nccurve[linecolor=red]{->}{l}{L}}
\end{frame}

\begin{frame}[fragile]{Working with strings}
\begin{block}<1->{Strings -- Synopsis}
\small
\begin{lstlisting}[language=Perl6,inputencoding=utf8,escapeinside={(*@}{@*)}]
my $s = chr(45) x 80; say (*@\pnode(0,0){A}{}@*)$s;
say "\c[WAVING WHITE FLAG, VARIATION SELECTOR-16,
        ZERO WIDTH JOINER, RAINBOW]".(*@\pnode(0,0){B}{}@*)ords;
my ($s1, $s2);
$s1 = 'Hello '; $s2 = "World\n";
print $s1 (*@\pnode(0,0){C}{}@*)~ $s2;
$s = "heLlo wOrld";
say "uc: ", $s.(*@\pnode(0,0){D}{}@*)uc, " -- tc: ", $s.(*@\pnode(0,0){E}{}@*)tc, " -- lc: ", $s.(*@\pnode(0,0){F}{}@*)lc;
\end{lstlisting}
\end{block}

\begin{itemize}
\item<2-> Outputs a textart line of \texttt{\rnode{a}{80 '-'}}
\item<3-> Tells all \rnode{b}{ordinal} values of the four codepoints of \texttt{\$s} (127987, 65039, 8205, 127752) % 🏳️‍🌈
\item<4-> \rnode{c}{Concatenates} the most important strings when learning a new language
\item<5-> \rnode{d}{UPPER} cases, \rnode{e}{Title} cases or \rnode{f}{lower} cases \lstinline[language=Perl6,inputencoding=latin9]!"heLlo wOrld"!
\end{itemize}
\uncover<2>{\nccurve[linecolor=teal]{->}{a}{A}}
\uncover<3>{\nccurve[linecolor=teal]{->}{b}{B}}
\uncover<4>{\nccurve[linecolor=teal]{->}{c}{C}}
\uncover<5>{\nccurve[linecolor=teal]{->}{d}{D}\nccurve[linecolor=teal]{->}{e}{E}\nccurve[linecolor=teal]{->}{f}{F}}
\end{frame}

\begin{frame}[fragile]{Match and Replacing}
\begin{block}<1->{Strings -- Synopsis}
\small
\begin{lstlisting}[language=Perl6,inputencoding=utf8,escapeinside={(*@}{@*)}]
my $r = "expletive expression";
$r ~~ s/expletive/(*@\pnode(0,0){A}{}@*)regular/;
my $s = "Hello world";
my $g = (*@\pnode(0,0){B}{}@*)~$/ given $s ~~ m/^(\w+)/;
(my $gpc = $s) ~~ (*@\pnode(0,0){C}{}@*)tr/oleH/terG/;
$gpc.substr-rw(*@\pnode(0,0){D}{}@*)(6,5) = "PerlCon";
\end{lstlisting}
\end{block}

\begin{itemize}
\item<2-> \rnode{a}{Text processing}!!! Not swearwords \ldots
\item<3-> What's our \rnode{b}{greeting} word?
\item<4-> \textit{Hello} \rnode{c}{\textrightarrow{}} \textit{Greet}
\item<5-> \textbf{Greet \rnode{d}{PerlCon}} 
\end{itemize}
\uncover<2>{\nccurve[linecolor=teal]{->}{a}{A}}
\uncover<3>{\nccurve[linecolor=teal]{->}{b}{B}}
\uncover<4>{\nccurve[linecolor=teal]{->}{c}{C}}
\uncover<5>{\nccurve[linecolor=teal]{->}{d}{D}}
\end{frame}


\subsection{Numbers}

\begin{frame}[fragile]{Int/Num Literals}
\begin{block}<1->{Numbers -- Synopsis}
\small
\begin{lstlisting}[language=Perl6,inputencoding=utf8,escapeinside={(*@}{@*)}]
say 10_(*@\pnode(0,0){B1}{}@*)00(*@\pnode(0,0){A1}{}@*)000 == 1_(*@\pnode(0,0){B2}{}@*)000(*@\pnode(0,0){A2}{}@*)_(*@\pnode(0,0){B3}{}@*)000;
say 1000(*@\pnode(0,0){A3}{}@*)000 - 1e(*@\pnode(0,0){A4}{}@*)6;
my $db = 0x(*@\pnode(0,0){C}{}@*)deadbeef;
chmod 0o(*@\pnode(0,0){D}{}@*)644, $?FILE;
say 0o755 - :(*@\pnode(0,0){E}{}@*)3<200021>;
say -(*@\pnode(0,0){F}{}@*)42;
\end{lstlisting}

\end{block}

\begin{itemize}
\item<2-> Four ways to note your first \rnode{a}{million} \uncover<3->{- mind the \rnode{b}{underscore} don't carry any semantic information}
\item<4-> Integer in \rnode{c}{hexadecimal} notation
\item<5-> \ldots and as typical in such a situation an \rnode{d}{octal} value
\item<6-> \rnode{e}{arbitary} bases can be used, too
\item<7-> \rnode{f}{unary operator} applied to literal \ldots but nitpicking
\end{itemize}
\uncover<2>{\nccurve[linecolor=teal]{->}{a}{A1}\nccurve[linecolor=teal]{->}{a}{A2}\nccurve[linecolor=teal]{->}{a}{A3}\nccurve[linecolor=teal]{->}{a}{A4}}
\uncover<3>{\nccurve[linecolor=teal]{->}{b}{B1}\nccurve[linecolor=teal]{->}{b}{B2}\nccurve[linecolor=teal]{->}{b}{B3}}
\uncover<4>{\nccurve[linecolor=teal]{->}{c}{C}}
\uncover<5>{\nccurve[linecolor=teal]{->}{d}{D}}
\uncover<6>{\nccurve[linecolor=teal]{->}{e}{E}}
\uncover<7>{\nccurve[linecolor=teal]{->}{f}{F}}
\end{frame}

\begin{frame}[fragile]{Complex Literals}
\begin{block}<1->{Numbers -- Synopsis}
\small
\begin{lstlisting}[language=Perl6,inputencoding=utf8,escapeinside={(*@}{@*)}]
say (*@\pnode(0,0){A}{}@*)i;
say (*@\pnode(0,0){B}{}@*)17i;
say (*@\pnode(0,0){C}{}@*)47+11i;
say -4+(*@\pnode(0,0){D}{}@*)Inf\i;
say 13-(*@\pnode(0,0){E}{}@*)8i;
say 2*(*@\pnode(0,0){F}{}@*)-3-2i;
say 2*(*@\pnode(0,0){G}{}@*)<-3-2i>;
\end{lstlisting}
\end{block}

\begin{itemize}
\item<2-> \rnode{a}{\texttt{0+1i}} - square root of \texttt{-1}
\item<3-> \rnode{b}{\texttt{0+17i}} - only the imaginary part
\item<4-> \rnode{c}{\texttt{47+11i}} - half real and half imaginary Genuine Eau de Cologne
\item<5-> \rnode{d}{\texttt{-4+Inf\textbackslash{}i}} - negative complex number with infinity imaginary part
\item<6-> \rnode{e}{\texttt{13-8i}} - positive complex number with negative imaginary part
\item<7-> \rnode{f}{\texttt{-6-2i}} - not a literal but mind operator precedence
\item<8-> \rnode{g}{\texttt{-6-4i}} - or better write complex number as literal
\end{itemize}
\uncover<2>{\nccurve[linecolor=teal]{->}{a}{A}}
\uncover<3>{\nccurve[linecolor=teal]{->}{b}{B}}
\uncover<4>{\nccurve[linecolor=teal]{->}{c}{C}}
\uncover<5>{\nccurve[linecolor=teal]{->}{d}{D}}
\uncover<6>{\nccurve[linecolor=teal]{->}{e}{E}}
\uncover<7>{\nccurve[linecolor=teal]{->}{f}{F}}
\uncover<8>{\nccurve[linecolor=teal]{->}{g}{G}}
\end{frame}

\begin{frame}[fragile]{Rational Literals}
\begin{block}<1->{Numbers -- Synopsis}
\small
\begin{lstlisting}[language=Perl6,inputencoding=utf8,escapeinside={(*@}{@*)}]
my $p = 3.1415(*@\pnode(0,0){A3}{}@*)9265; $p = pi(*@\pnode(0,0){A1}{}@*); $p = π(*@\pnode(0,0){A2}{}@*);
my $t = τ(*@\pnode(0,0){B}{}@*);
my $e = e(*@\pnode(0,0){C}{}@*);
say ½(*@\pnode(0,0){D}{}@*).perl, ", ", <1/(*@\pnode(0,0){E}{}@*)7>.perl;
\end{lstlisting}

\end{block}

\begin{itemize}
\item<2-> Calculate \rnode{a1}{\textit{2 * \textpi{} * r}} or \rnode{a2}{\textit{\textpi{} * r\textsuperscript{2}}} using the correct constant instead of hard-coded \rnode{a3}{NIH} value
\item<3-> \rnode{b}{\textit{τ=2π}}
\item<4-> \rnode{c}{Euler's} number \ldots
\item<5-> Outputs \rnode{d}{\texttt{0.5}}\texttt{, }\rnode{e}{\texttt{<1/7>}}
\end{itemize}
\uncover<2>{\nccurve[linecolor=teal]{->}{a1}{A1}\nccurve[linecolor=teal]{->}{a2}{A2}\nccurve[linecolor=teal]{->}{a3}{A3}}
\uncover<3>{\nccurve[linecolor=teal]{->}{b}{B}}
\uncover<4>{\nccurve[linecolor=teal]{->}{c}{C}}
\uncover<5>{\nccurve[linecolor=teal]{->}{d}{D}\nccurve[linecolor=teal]{->}{e}{E}}
\end{frame}

\begin{frame}[fragile]{Basic operations}
\begin{block}<1->{Numbers -- Synopsis}
\small
\begin{lstlisting}[language=Perl6,inputencoding=utf8,escapeinside={(*@}{@*)}]
say $t <=> $e; # (*@\pnode(0,0){A}{}@*)<, <=, ==, >=, >, !=, <=>
say π+e;       # +, -, *, /, (*@\pnode(0,0){B}{}@*)%, **, %%
say $p++, " -- ", ++$p;      # (*@\pnode(0,0){C}{}@*)++, --, -, +, ^, ?, ~
say 1 +< 4;    # +<, +>, (*@\pnode(0,0){D}{}@*)+&, +|
say π(*@\pnode(0,0){E}{}@*)²;        # ¹²³⁴⁵⁶⁷⁸⁹⁰
say 2¹²(*@\pnode(0,0){F}{}@*)⁷-1;    # all directly appended superscripted belong together like any number
\end{lstlisting}

\end{block}

\begin{itemize}
\item<2-> \rnode{a}{comparison} operators
\item<3-> binary \rnode{b}{arithmetic} operators
\item<4-> \rnode{c}{prefix} -- postfix operators \ldots and other unary operators
\item<5-> shift left as shortcut for power of 2 - binary \rnode{d}{bitwise} operators
\item<6-> direct power of \rnode{e}{arbitary} base
\item<7-> can be \rnode{f}combined to a very large prime number
\end{itemize}
\uncover<2>{\nccurve[linecolor=teal]{->}{a}{A}}
\uncover<3>{\nccurve[linecolor=teal]{->}{b}{B}}
\uncover<4>{\nccurve[linecolor=teal]{->}{c}{C}}
\uncover<5>{\nccurve[linecolor=teal]{->}{d}{D}}
\uncover<6>{\nccurve[linecolor=teal]{->}{e}{E}}
\uncover<7>{\nccurve[linecolor=teal]{->}{f}{F}}
\end{frame}


\subsection{Conditionals}

\begin{frame}[fragile]{Conditionals - If}
\begin{block}<1->{If -- Synopsis}
\small
\begin{lstlisting}[language=Perl6,inputencoding=utf8,escapeinside={(*@}{@*)}]
i(*@\pnode(0,0){A}{}@*)f True { say "Yes" } el(*@\pnode(0,0){B}{}@*)se { say "No" };
if incisors and swimming (*@\pnode(0,0){C}{}@*){
    i(*@\pnode(0,0){D}{}@*)f    broad_tail     { say "Beaver (Castor)" }
    elsi(*@\pnode(0,0){E}{}@*)f flattened_tail { say "Muskrat" }
    els(*@\pnode(0,0){F}{}@*)e  small_tail     { say "Coypu" }
}
say "\c[LIFE WITHOUT ALFA ROMEO]" un(*@\pnode(0,0){G}{}@*)less Alfa-Driver;
say "\c[LIFE WITH ALFA ROMEO]" i(*@\pnode(0,0){H}{}@*)f Alfa-Driver;
\end{lstlisting}
\end{block}

\begin{itemize}
\item<2-> A very ordinary \rnode{a}{\texttt{if}} \ldots \rnode{b}{\texttt{else}} line
\item<3-> A larger \rnode{c}{\texttt{if}}-block containing some \rnode{d}{\texttt{if}}, \rnode{e}{\texttt{elsif}} (not, but \ldots) and finally an \rnode{f}{\texttt{else}}
\item<4-> statement modifier printing \includegraphics[height=\baselineskip]{Life-without-Alfa-Romeo.eps} for \rnode{g}{unfortunate} ones (currently missing in unicode)
\item<5-> statement modifier printing \includegraphics[height=\baselineskip]{Life-with-Alfa-Romeo.eps} for \rnode{h}{fortunate} ones (currently missing in unicode, too)
\end{itemize}
\uncover<2>{\nccurve[linecolor=teal]{->}{a}{A}\nccurve[linecolor=teal]{->}{b}{B}}
\uncover<3>{\nccurve[linecolor=teal]{->}{c}{C}\nccurve[linecolor=teal]{->}{d}{D}\nccurve[linecolor=teal]{->}{e}{E}\nccurve[linecolor=teal]{->}{f}{F}}
\uncover<4>{\nccurve[linecolor=teal]{->}{g}{G}}
\uncover<5>{\nccurve[linecolor=teal]{->}{h}{H}}
\end{frame}

\begin{frame}[fragile]{Conditionals - With}
\begin{block}<1->{With -- Synopsis}
\small
\begin{lstlisting}[language=Perl6,inputencoding=utf8,escapeinside={(*@}{@*)}]
with "abcdefg".index(("a".."g").roll) { .s(*@\pnode(0,0){A}{}@*)ay }
without (Any, True).(*@\pnode(0,0){B1}{}@*)roll { say "Roll (*@\pnode(0,0){B2}{}@*)again" }
sub prime-only { my $v = (^2019).pick;
  $v.is-prime || return Any; return $v }
with prime-only() -> (*@\pnode(0,0){C}{}@*)$n { say "$n is prime" }
my $s = ("a".."g").roll;
with(*@\pnode(0,0){D}{}@*)   $s.index("a") { say "say <b> now" }
o(*@\pnode(0,0){E}{}@*)rwith $s.index("b") { say "did you say <a>?" }
else   { say "neither <a> nor <b> was said" }
\end{lstlisting}
\end{block}

\begin{itemize}
\item<2-> \rnode{a}{says} a number from \texttt{0 \ldots 6}
\item<3-> \rnode{b1}{50\%} chance to get \rnode{b2}{\textit{Roll again}}
\item<4-> way lower chance to see \textit{\ldots is prime} using \rnode{c}{pointy block} syntax together with \texttt{with}
\item<5-> analogous to \texttt{if \ldots elsif \ldots} you can concatenate multiple tests using \texttt{\rnode{e}{o}rwith} after \rnode{d}{\texttt{with}}
\end{itemize}
\uncover<2>{\nccurve[linecolor=teal]{->}{a}{A}}
\uncover<3>{\nccurve[linecolor=teal]{->}{b1}{B1}\nccurve[linecolor=teal]{->}{b2}{B2}}
\uncover<4>{\nccurve[linecolor=teal]{->}{c}{C}}
\uncover<5>{\nccurve[linecolor=teal]{->}{d}{D}\nccurve[linecolor=teal]{->}{e}{E}}
\end{frame}

\begin{frame}[fragile]{Conditionals - Given/When}
\begin{block}<1->{Given -- Synopsis}
\small
\begin{lstlisting}[language=Perl6,inputencoding=utf8,escapeinside={(*@}{@*)}]
given  (2, 3, 7, "Any", Any).pick {
  when 2     { say $_ * 3 * 7; }(*@\pnode(0,0){A}{}@*)
  when 3     { say $_ * 2 * 7; }(*@\pnode(0,0){B}{}@*)
  when 7     { say $_ * 2 * 3; }(*@\pnode(0,0){C}{}@*)
  when "Any" { say "Any" }(*@\pnode(0,0){D}{}@*)
  default    { say "(*@\pnode(0,0){E}{}@*)Hmpf"; }
}
\end{lstlisting}
\end{block}

\begin{itemize}
\item<2-> 60\% chance to get the answer to the ultimate question of life, the universe, and everythin\rnode{a}{g}
\item<3-> 20\% chance to get \textit{\rnode{d}{Any}} answer
\item<4-> 20\% chance to \rnode{e}{disappoint} the automate
\end{itemize}
\uncover<2>{\nccurve[linecolor=teal]{->}{a}{A}\nccurve[linecolor=teal]{->}{a}{B}\nccurve[linecolor=teal]{->}{a}{C}}
\uncover<3>{\nccurve[linecolor=teal]{->}{d}{D}}
\uncover<4>{\nccurve[linecolor=teal]{->}{e}{E}}
\end{frame}

\begin{frame}[fragile]{Conditionals - Given/When}
\begin{block}<1->{Given -- Synopsis}
\small
\begin{lstlisting}[language=Perl6,inputencoding=utf8,escapeinside={(*@}{@*)}]
given  (2, 3, 7, "Any", Any).pick {
  when Int   (*@\pnode(0,0){A}{}@*){ say "Int found, computing ..."; }
  when 2     { say $_ * 3 * 7; }(*@\pnode(0,0){B1}{}@*)
  when 3     { say $_ * 2 * 7; }(*@\pnode(0,0){B2}{}@*)
  when 7     { say $_ * 2 * 3; }(*@\pnode(0,0){B3}{}@*)
  when Str   (*@\pnode(0,0){C}{}@*){ say "Str found, checking ..."; }
  when "Any" { say "Any" }(*@\pnode(0,0){D}{}@*)
  default    { say "(*@\pnode(0,0){E}{}@*)Hmpf"; }
}
\end{lstlisting}
\end{block}

\begin{itemize}
\item<1-> The block surrounding \texttt{when} or \texttt{default} statements will be left immediately after leaving the sub-block after \texttt{when} or \texttt{default} statement.
\item<2-> Read: \rnode{a}{1st} \texttt{when} block skips \rnode{b1}{2nd}, \rnode{b2}{3rd} and \rnode{b3}{4th} \texttt{when} block\uncover<3->{,\rnode{c}{5th} \texttt{when} block skips \rnode{d}{6th} \texttt{when} block}
\item<4-> \texttt{when} statement modifiers won't cause leaving a block
\end{itemize}
\uncover<2>{\nccurve[linecolor=teal]{->}{a}{A}\nccurve[linecolor=teal]{->}{b1}{B1}\nccurve[linecolor=teal]{->}{b2}{B2}\nccurve[linecolor=teal]{->}{b3}{B3}}
\uncover<3>{\nccurve[linecolor=teal]{->}{c}{C}\nccurve[linecolor=teal]{->}{d}{D}}
\end{frame}

\begin{frame}[fragile]{Conditionals - Proceed/Succeed}
\begin{block}<1->{Given -- Synopsis}
\small
\begin{lstlisting}[language=Perl6,inputencoding=utf8,escapeinside={(*@}{@*)}]
given  (2, 3, 7, "Any", Any).pick {
  when Int   { print "Int found, computing ..."; proc(*@\pnode(0,0){A}{}@*)eed; }
  when 2     { say $_ * 3 * 7; }(*@\pnode(0,0){B1}{}@*)
  when 3     { say $_ * 2 * 7; }(*@\pnode(0,0){B2}{}@*)
  when 7     { say $_ * 2 * 3; }(*@\pnode(0,0){B3}{}@*)
  when Str   { say "Str found, checking ..."; proc(*@\pnode(0,0){C}{}@*)eed; }
  when "Any" { say "Any" }(*@\pnode(0,0){D}{}@*)
  default    { su(*@\pnode(0,0){E}{}@*)cceed; say "Hmpf"; }
}
\end{lstlisting}
\end{block}

\begin{itemize}
\item<2->  \rnode{a}{Now} there is a 60\% chance to get the answer to the ultimate question of life, the universe, and everything \textbf{computed}
\item<3-> 20\% to get \rnode{c}{now} a checked \textit{Any} - \texttt{Str} as output
\item<4-> In \texttt{Any} other case (\texttt{default}), the automate \rnode{e}{immediately succeeds} and keeps its disappointment for itself
\end{itemize}
\uncover<2>{\nccurve[linecolor=teal]{->}{a}{A}}
\uncover<3>{\nccurve[linecolor=teal]{->}{c}{C}}
\uncover<4>{\nccurve[linecolor=teal]{->}{e}{E}}
\end{frame}

\subsection{Conditionals}

\begin{frame}[fragile]{Loops}
\begin{block}<1->{for -- Synopsis}
\small
\begin{lstlisting}[language=Perl6,inputencoding=utf8,escapeinside={(*@}{@*)}]
my @l = (^10).pick(3);
for(*@\pnode(0,0){A}{}@*) @l { .(*@\pnode(0,0){B}{}@*)say }
for @l <(*@\pnode(0,0){C}{}@*)-> $n { ++$n }
for @l ->(*@\pnode(0,0){D}{}@*) $a, $b = "n/a" { say "a: $a, b: $b"; }
@l = (^10).pick(4).sort;
for @l { say "(*@\pnode(0,0){E}{}@*)$^y > $^x" }
\end{lstlisting}
\end{block}

\begin{itemize}
\item<2-> \rnode{a}{loop} over elemenst in \texttt{@l}, put the current one into \texttt{\$\_} and execute the block running \rnode{b}{\texttt{\$\_.say}}
\item<3-> run a similar loop and write results \rnode{c}{back}
\item<4-> just another \rnode{d}{pointy block} with default value \ldots -- outputs something like \texttt{\\
2, 7 \\
8, n/a}
\item<5-> loops over \texttt{@l} taking \rnode{e}{two elements} at once and put them into \texttt{(\$x, \$y)} -- alphabetically ordered
\end{itemize}
\uncover<2>{\nccurve[linecolor=teal]{->}{a}{A}\nccurve[linecolor=teal]{->}{b}{B}}
\uncover<3>{\nccurve[linecolor=teal]{->}{c}{C}}
\uncover<4>{\nccurve[linecolor=teal]{->}{d}{D}}
\uncover<5>{\nccurve[linecolor=teal]{->}{e}{E}}
\end{frame}

\begin{frame}[fragile]{Loops}
\begin{block}<1->{loop -- Synopsis}
\small
\begin{lstlisting}[language=Perl6,inputencoding=utf8,escapeinside={(*@}{@*)}]
(*@\pnode(0,0){A}{}@*)loop (my $i = 0; $i < 10; ++$i) { $i.say }
loop (*@\pnode(0,0){B}{}@*){ say 'forever' }
my @l =(*@\pnode(0,0){C}{}@*) loop (my $i = 10; $i < 20; ++$i) { $i + $i - 1 }
say @l;
(loop (my $i = 10; $i < 20; ++$i) { $i * 2 - 1 }).(*@\pnode(0,0){D}{}@*)say;
\end{lstlisting}%$
\end{block}

\begin{itemize}
\item<2-> \rnode{a}{loop} as known from C programming language
\item<3-> easy to write \rnode{b}{infinite} loops - no need for hacks
\item<4-> loops can be used to produce \rnode{c}{results} \ldots
\item<5-> \ldots which can be processed \rnode{d}{immediately}
\end{itemize}
\uncover<2>{\nccurve[linecolor=teal]{->}{a}{A}}
\uncover<3>{\nccurve[linecolor=teal]{->}{b}{B}}
\uncover<4>{\nccurve[linecolor=teal]{->}{c}{C}}
\uncover<5>{\nccurve[linecolor=teal]{->}{d}{D}}
\end{frame}

\begin{frame}[fragile]{Loops}
\begin{block}<1->{while -- Synopsis}
\small
\begin{lstlisting}[language=Perl6,inputencoding=utf8,escapeinside={(*@}{@*)}]
my $i = 1; my $p = 1;
while $i <(*@\pnode(0,0){A}{}@*) 5 {(*@\pnode(0,0){B1}{}@*) ++$i if (++$p).is-prime (*@\pnode(0,0){B2}{}@*)}
say "$p is the {$i}th prime";
my $x = 1;
"{$x**2} <= {2**$x}".say and ++$x whi(*@\pnode(0,0){C}{}@*)le $x**2 <= 2**$x;
\end{lstlisting}%$
\end{block}

\begin{itemize}
\item<2-> \texttt{while} runs the loop as long as the \rnode{a}{\textit{expression}} evaluates to \texttt{true}
\item<3-> \textit{expression} is evaluated before the block is executed\uncover<4->{ -- might result in entire \rnode{b}{block} is skipped}
\item<5-> \texttt{while} can be used as statement \rnode{c}{modifier}, too
\end{itemize}
\uncover<2>{\nccurve[linecolor=teal]{->}{a}{A}}
\uncover<4>{\nccurve[linecolor=teal]{->}{b}{B1}\nccurve[linecolor=teal]{->}{b}{B2}}
\uncover<5>{\nccurve[linecolor=teal]{->}{c}{C}}
\end{frame}

\begin{frame}[fragile]{Loops}
\begin{block}<1->{until -- Synopsis}
\small
\begin{lstlisting}[language=Perl6,inputencoding=utf8,escapeinside={(*@}{@*)}]
my $i = 1; my $p = 1;
until $i (*@\pnode(0,0){A}{}@*)> 5 {(*@\pnode(0,0){B1}{}@*) ++$i if (++$p).is-prime (*@\pnode(0,0){B2}{}@*)}
say "$p is the {$i}th prime";
my $x = 1;
$x++ un(*@\pnode(0,0){C}{}@*)til "$x".chars > 2;
\end{lstlisting}%$
\end{block}

\begin{itemize}
\item<2-> \texttt{until} runs the loop as long as the \rnode{a}{\textit{expression}} evaluates to \texttt{false}
\item<3-> \textit{expression} is evaluated before the block is executed\uncover<4->{ -- might result in entire \rnode{b}{block} is skipped}
\item<5-> \texttt{until} can be used as statement \rnode{c}{modifier}, too
\end{itemize}
\uncover<2>{\nccurve[linecolor=teal]{->}{a}{A}}
\uncover<4>{\nccurve[linecolor=teal]{->}{b}{B1}\nccurve[linecolor=teal]{->}{b}{B2}}
\uncover<5>{\nccurve[linecolor=teal]{->}{c}{C}}
\end{frame}

\begin{frame}[fragile]{Loops}
\begin{block}<1->{repeat \ldots while/until -- Synopsis}
\small
\begin{lstlisting}[language=Perl6,inputencoding=utf8,escapeinside={(*@}{@*)}]
my $i = 0;
repeat { $i++ } while $i (*@\pnode(0,0){A}{}@*)< 10;
say $i;
repeat while $i (*@\pnode(0,0){B}{}@*)< 10 { $i++ };
say $i;
repeat { $i++ } until "$i".chars (*@\pnode(0,0){C}{}@*)> 2;
say $i;
\end{lstlisting}
\end{block}

\begin{itemize}
\item<2-> \texttt{repeat} runs the block at least once -- \rnode{a}{\textit{expression}} is evaluated after the block has been run
\item<3-> \texttt{repeat \ldots while} runs as long as the expression evaluates to \rnode{b}{\texttt{true}}
\item<4-> \texttt{repeat \ldots until} runs as long as the expression evaluates to \rnode{c}{\texttt{false}}
\end{itemize}
\uncover<2>{\nccurve[linecolor=teal]{->}{a}{A}}
\uncover<3>{\nccurve[linecolor=teal]{->}{b}{B}}
\uncover<4>{\nccurve[linecolor=teal]{->}{c}{C}}
\end{frame}




%\begin{frame}[fragile]{Literals}
%\begin{block}<1->{Numbers -- Synopsis}
%\small
%\begin{lstlisting}[language=Perl6,inputencoding=utf8,escapeinside={(*@}{@*)}]
%my $p = 3.14159265; $p = pi; $p = π;
%\end{lstlisting}
%\end{block}
%
%\begin{itemize}
%\item<2-> \rnode{a}{}
%\item<3-> 
%\item<4-> 
%\item<5-> 
%\end{itemize}
%\uncover<2>{\nccurve[linecolor=teal]{->}{a}{A}}
%\uncover<3>{\nccurve[linecolor=teal]{->}{b}{B}}
%\uncover<4>{\nccurve[linecolor=teal]{->}{c}{C}}
%\uncover<5>{\nccurve[linecolor=teal]{->}{d}{D}}
%\end{frame}

\section{Thank you}

% THANKS
\begin{frame}[fragile]{Thank You For Listening}
\begin{block}<1->{Questions?}
Jens Rehsack \textless{}\href{mailto:rehsack@cpan.org}{rehsack@cpan.org}\textgreater{} \\
Cologne
\end{block}
\end{frame}

\end{document}

